\documentclass[10pt,a4paper]{article}
\usepackage[margin=1.75cm]{geometry}
\usepackage[T1]{fontenc}
\usepackage{lmodern}
\usepackage{microtype}
\usepackage{booktabs}
\usepackage{tabularx}
\usepackage{array}
\usepackage{longtable}
\usepackage{hyperref}
\hypersetup{colorlinks=true,urlcolor=blue,linkcolor=black}
\setlength{\parskip}{0.45em}
\setlength{\parindent}{0pt}
\newcolumntype{Y}{>{\raggedright\arraybackslash}X}
\newcommand{\obs}{\textbf{Observed.} }
\newcommand{\inferr}{\textbf{Inference.} }
\newcommand{\caveat}{\textbf{Caveat.} }
\title{RelaLeap: What Was Actually Tested in the Predictive-Coding / ePC Track?}
\author{Consolidated experiment report\\RelaLeap project record}
\date{6 August 2026}
\begin{document}
\maketitle

\section*{Executive summary}
This report corrects an important naming issue in the earlier experiment history.
The runs originally called ``ePC'' (especially R6--R9) used a \emph{state-based}
predictive-coding implementation: hidden states were optimized and then detached
before parameter differentiation.  A later source and method audit established
that this is SO/sPC, not the error-variable EO/ePC method described in the
reference paper.  Therefore their negative performance results are evidence
against that \emph{legacy implementation}, not a clean test or refutation of
genuine EO/ePC.

The legacy programme did produce one robust positive engineering observation:
its declared activity-energy decreased monotonically in 10/10 multi-seed
records.  But it did not turn that property into practical advantages.  On a
six-layer GPT-2 student trained on WikiText-103, ordinary backpropagation with
knowledge distillation (KD), given the same elapsed time, made about 6.7 times
as many updates and had substantially better validation loss.  Legacy PC also
showed lower representation rank and weaker causal-factor selectivity.

The subsequent four-regime CPU preflight likewise does \emph{not} establish an
ePC efficacy result: its ordinary T=1 control could not learn the fixture.
Moreover, the later audit found a metric defect that selected the worst rather
than the best cap.  Its old arm comparisons are invalid for method comparison.

The replacement EO/ePC implementation has passed local mechanism and
instrument checks: it has direct error-variable dynamics, settles to the exact
PC fixed point in small analytical fixtures, has a validated target/evaluation
seam, and replays bit-exactly.  This is validity evidence, not yet a
language-model efficacy result.  The planned five-seed GPT-2/WikiText reference
reproduction (C1) has not successfully run; it remains blocked pending a fresh,
single-writer GPU execution envelope.

\section{How to read the evidence}
\begin{itemize}
\item \emph{Observed} means an artifact or recorded result directly supports the statement.
\item \emph{Inference} distinguishes what the results suggest from what they prove.
\item A completed software or analytic control is not treated as evidence that EO/ePC
  improves language-model learning.
\end{itemize}

\section{Chronology and status}
\begin{longtable}{p{0.13\linewidth}p{0.23\linewidth}p{0.27\linewidth}p{0.28\linewidth}}
\toprule
Date & Run & Question & Result / status \\
\midrule
\endfirsthead
\toprule Date & Run & Question & Result / status \\
\midrule
\endhead
19 Jul & R6 mechanism screen & Can five arms complete and expose dynamics? & Failed after the two BP arms; no PC result. \\
20 Jul & R7 one-seed screen & Does the repaired five-arm legacy run complete? & Completed; small loss edge over update-matched KD, no benefit from deeper inference. Exploratory only. \\
20--21 Jul & R8, five seeds & Does legacy PC beat BP/KD on GPT-2/WikiText? & No: BP-KD at matched time dominated; energy monotonicity held for legacy PC. \\
21--22 Jul & R9 representation battery & Does legacy PC yield better internal representations? & No: lower rank and lower causal selectivity than BP-CE. \\
4 Aug & Multiregime preflight & Is a larger planted-task fixture learnable before a frontier/MoE study? & No: the matched T=1 control failed every regime; no frontier study was admitted. Later audit also invalidated its old comparison metric. \\
5 Aug onward & EO/ePC remediation & Is the method and measurement apparatus actually EO/ePC and trustworthy? & Local method, fixed-point, target-seam, and evaluator controls passed. No language efficacy result yet. \\
\bottomrule
\end{longtable}

\section{Legacy experiments (R6--R9)}
\subsection{R6: interrupted mechanism screen}
R6 was a one-seed, 200-update GPU screen on a six-layer GPT-2 student and
WikiText-103.  It compared BP cross-entropy (CE), update-matched BP-KD, two
legacy PC depths (T=4 and T=12), and wall-clock-matched BP-KD.  The job failed
after CE and KD completed, before any legacy-PC arm completed.  Its immediate
exception was not retained.  Thus it supports no PC conclusion.  The two saved
baselines finished at validation losses 8.439 (CE) and 7.628 (KD).

\subsection{R7: repaired exploratory screen}
R7 fixed the arm-loop crash and completed all five arms for one seed.  At 200
updates the reported validation losses were: CE 8.4390, KD 7.6278, legacy PC
T=4 7.5672, legacy PC T=12 7.5836, and wall-clock KD 7.6278.  This was an
interesting single-seed hint, but it was not a confidence interval or a
decisive comparison.  Increasing inference depth fourfold did not help.

\subsection{R8: five-seed outcome comparison}
R8 repeated the design across five frozen seeds (25 records).  It is the main
legacy performance result.

\begin{center}
\small
\begin{tabular}{lrrrrr}
\toprule
Arm & mean validation loss & mean perplexity & updates & elapsed seconds & interpretation \\
\midrule
BP-CE & $6.261\pm0.060$ & 524 & 200 & 78 & best equal-update baseline \\
BP-KD & $7.039\pm0.056$ & 1142 & 200 & 74 & KD noisy at 200 updates \\
Legacy PC T=4 & $6.781\pm0.146$ & 889 & 200 & 445 & slower, worse than CE \\
Legacy PC T=12 & $6.754\pm0.131$ & 863 & 200 & 1738 & no convincing depth gain \\
BP-KD, matched time & $5.785\pm0.135$ & 328 & 1334 & 445 & dominant matched-time arm \\
\bottomrule
\end{tabular}
\end{center}

\obs All 10 legacy-PC seed/arm records satisfied the declared
activity-energy-monotonicity check.  In the time taken by T=4 legacy PC to
make 200 updates, ordinary KD made 1334: approximately 6.7 times more.

\inferr The evidence says that the legacy state-PC implementation did not
provide a useful speed/quality trade-off on this task.  It does not establish
that EO/ePC would behave identically, because the algorithm was not the same.

\subsection{R9: representation and causal-factor diagnostics}
R9 analyzed 20 of the R8 trained models.  At the final layer, effective rank
was $13.9\pm0.3$ for BP-CE, $5.0\pm0.1$ for BP-KD, $3.9\pm0.2$ for legacy PC
T=4, and $3.8\pm0.1$ for legacy PC T=12.  Participation ratio was 6.1 for
CE, 3.1 for KD, and about 2.6 for both legacy-PC arms.  These are strong signs
of representational compression, not dead units (dead fraction was 0.001).

Final-layer causal-selectivity probes also favored BP-CE.  For subject number,
object number, tense, and negation respectively, BP-CE scored 0.607, 0.425,
0.661, and 0.870; legacy T=4 scored 0.438, 0.232, 0.546, and 0.719; legacy
T=12 scored 0.379, 0.275, 0.481, and 0.718.  The original interpretation was
that smoother activity energy coincided with more compressed but less
factor-selective representations.  That remains a reasonable description of
these legacy runs, but not a general theorem about EO/ePC.

\section{The multiregime frontier preflight}
The four-regime CPU preflight was meant to admit a richer planted transformer
task before testing a Pareto frontier of PC variants or an MoE-like conditional
specialization scheme.  Twelve arm-by-seed executions replayed exactly; T=2 and
T=4 were energy-monotone.  But the required T=1 control had NLLs far above its
1.4 learnability floor in every regime (roughly 6.76--11.92).  The fixture was
therefore correctly stopped before a frontier claim.

\caveat A later source audit established both that legacy code was SO/sPC and
that a negation in the old cap-selection path selected the worst, rather than
best, cap.  Hence the values from this preflight cannot be used to compare
methods.  What survives is the narrow operational fact: the fixture, as then
specified, failed the ordinary-control learnability gate.

\section{EO/ePC remediation: what has now been established}
After the method audit, the track was rebuilt around the paper's error-variable
computation.  The following are local, reproducible validation steps.

\begin{tabularx}{\linewidth}{p{0.15\linewidth}Y Y}
\toprule
Gate & Observed result & Meaning and limit \\
\midrule
A0 method/instrument audit & Pinned paper and official-code sources matched; uniform NLL was $\ln(64)$; planted-oracle NLL was about $2.38\times10^{-9}$; trainer-teacher round trip was 0.0. & Established old routine was SO/sPC and validated basic measurement paths. \\
MG-1 parameterization & At depths 2/6/12/24, every first error update equaled the output adjoint times -0.01 (relative error 0). Old state-based path had zero first-step update in earlier layers. & Establishes direct deep error signal in the new EO/ePC core; no efficacy claim. \\
MG-2R fixed-point check & All 12 EO points settled (relative residual $5.53\times10^{-9}$ to $8.14\times10^{-9}$), matched exact controls, and replayed bit-exactly. & EO/ePC realizes the finite-lambda PC fixed point in the tested analytical fixtures. \\
MG-3/MG-4 settling/stability & Explicit residual certificates and local map-stability estimates passed their frozen controls. & Shows numerical diagnostics behave as declared, only for the tested maps. \\
MG-5/MG-6 trainer seam & All 512 planted target positions had a unique correct target; corruptions failed closed; evaluation/teacher round trips passed. & Target production and scoring are no longer the known source of a false conclusion. \\
\bottomrule
\end{tabularx}

One earlier MG-2 formulation failed because it demanded finite-lambda equality
to BP at tolerance $10^{-6}$, which a seed-free scalar counterexample showed
to be mathematically impossible.  MG-2R replaced that invalid demand with the
correct lambda-to-zero convergence-rate criterion and passed.  This is a repair
of a test contract, not post-hoc tuning of an efficacy outcome.

\section{What is known now, and what is not}
\textbf{Known.} The old state-PC experiments were completed and their raw
outcomes are preserved.  They show monotone declared activity energy but no
competitive learning or representation benefit against their baselines.
The new core is materially different: it has passed analytic EO/ePC mechanism
and measurement controls.

\textbf{Not known.} Whether genuine EO/ePC improves a transformer on
language-model learning, distillation, robustness, adaptation, or causal
representation quality is still untested.  The relevant next study is the
frozen C1 five-seed reproduction with GPT-2 and WikiText.  Its local package
and adjudication seams have been built, but no valid C1 science run has
completed; prior remote attempts were terminated before model/data/seed work
because of package and provider-control failures.

\section*{Primary evidence and reproducibility}
The project ledgers are authoritative.  This report is a synthesis, not a
replacement for raw artifacts.
\begin{itemize}
\item \texttt{experiments/20260719T202100Z-epc-mechanism-screen-r6/RUN.md}
\item \texttt{experiments/20260720T114200Z-epc-mechanism-screen-r7/RUN.md}
\item \texttt{experiments/20260720T193211Z-epc-r8-broad-outcome/RUN.md}
\item \texttt{experiments/20260721T222400Z-r9-stage-a/RUN.md}
\item \texttt{experiments/20260804T173241Z-epc-multiregime-frontier-preflight-v1-r5/RUN.md}
\item \texttt{experiments/20260805T002848Z-epc-remediation-a0-acceptance-r2/RUN.md}
\item \texttt{experiments/20260805T004556Z-epc-remediation-mg1-artifact-acceptance/RUN.md}
\item \texttt{experiments/20260805T015749Z-epc-remediation-mg2r-acceptance/RUN.md}
\item \texttt{experiments/20260805T035047Z-epc-remediation-mg5-acceptance/RUN.md}
\end{itemize}
The earlier historical companion is
\texttt{docs/prior\_epc\_experiment\_history\_20260804.pdf}; this report
supersedes its method-identity interpretation where the later audit found the
SO/sPC versus EO/ePC distinction.

\end{document}
